\begin{mistake}{2026-04-12}{30}

\begin{problemcontent}
求微分方程 \( y''-y=x^2 \) 的解 \( \varphi(x) \)，使得 \( \varphi(x) \) 在 \( x\to 0 \) 时是较 \( x^2 \) 高阶的无穷小。
\end{problemcontent}

\begin{solutioncontent}
先求原非齐次方程的通解。

对应齐次方程为
\[
y''-y=0
\]
其特征方程为
\[
r^2-1=0
\]
故齐次通解为
\[
y_h=C_1e^x+C_2e^{-x}
\]

对非齐次项 \(x^2\)，设多项式特解为
\[
y_p=ax^2+bx+d
\]
则
\[
y_p''=2a
\]
代入方程 \(y''-y=x^2\)，得
\[
2a-(ax^2+bx+d)=x^2
\]
比较同次幂系数：
\[
-a=1,\qquad -b=0,\qquad 2a-d=0
\]
解得
\[
a=-1,\qquad b=0,\qquad d=-2
\]
故一个特解为
\[
y_p=-x^2-2
\]
于是原方程通解为
\[
y=C_1e^x+C_2e^{-x}-x^2-2
\]

现要求 \(\varphi(x)=o(x^2)\)（即比 \(x^2\) 高阶），故需将其在 \(x=0\) 处展开到二次项：
\[
e^x=1+x+\frac{x^2}{2}+o(x^2),\qquad e^{-x}=1-x+\frac{x^2}{2}+o(x^2)
\]
代入得
\[
\varphi(x)=C_1\left(1+x+\frac{x^2}{2}\right)+C_2\left(1-x+\frac{x^2}{2}\right)-x^2-2+o(x^2)
\]
整理为
\[
\varphi(x)=(C_1+C_2-2)+(C_1-C_2)x+\left(\frac{C_1+C_2}{2}-1\right)x^2+o(x^2)
\]
要使 \(\varphi(x)=o(x^2)\)，必须使常数项、一次项、二次项系数全为 0，即
\[
C_1+C_2=2,\qquad C_1-C_2=0
\]
解得
\[
C_1=C_2=1
\]
所以所求解为
\[
\varphi(x)=e^x+e^{-x}-x^2-2
\]
\end{solutioncontent}

\mistakeknowledge{
\begin{itemize}
  \item \textbf{核心公式}：常系数非齐次方程通解等于“齐次通解 + 一个特解”；比较系数法求多项式特解；\(o(x^2)\) 要看展开后到二次项是否全部消失。
  \item \textbf{易错点}：特解不能只设 \(ax^2\)，而应设完整二次式 \(ax^2+bx+d\)；“较 \(x^2\) 高阶”不是只消去常数项和一次项，还必须使二次项也为 0。
  \item \textbf{关联知识}：判断无穷小阶数时，常把通解在指定点作泰勒展开，再通过选常数消去低阶项，这是常见题型。
\end{itemize}}

\end{mistake}
